\documentclass[11pt]{article}
\usepackage[margin=1in]{geometry}
\usepackage{amsmath,amssymb,amsthm}
\usepackage{graphicx}
\usepackage[hidelinks]{hyperref}

\newtheorem{theorem}{Theorem}
\newtheorem{lemma}{Lemma}

\newcommand{\Lop}{\mathcal L}
\newcommand{\I}{I}

\title{A Separable Banach-Space Counterexample to Strong Approximation by Two Idempotents}
\author{fa\_banach\_001, agent\_lane\_01}
\date{August 26, 2026}

\begin{document}
\maketitle

\section*{Status}

Candidate counterexample, likely valid, pending human review.

\section*{Source Question}

Jean-Christophe Bourin, \emph{A Journey into Matrix Analysis},
arXiv:2307.03064, Question 5.2.9 on PDF page 82, asks:
\begin{quote}
Let $\mathcal X$ be a separable Banach space and
$T\in\Lop(\mathcal X)$. Do there exist two sequences
$\{P_n\}_{n=1}^{\infty}$ and $\{Q_n\}_{n=1}^{\infty}$ of idempotents in
$\Lop(\mathcal X)$ such that
\[
T=\operatorname{sot}\!\lim_{n\to\infty}(P_n+Q_n)?
\]
\end{quote}

\begin{center}
\includegraphics[width=.97\linewidth]{figures/open_problem_crop.png}
\end{center}

Here strong-operator convergence means pointwise norm convergence on
$\mathcal X$.

\section*{Idea of the Proof}

The Hilbert-space construction behind the source question uses the abundance of
large complemented subspaces.  The obstruction is to choose a separable space at
the opposite extreme.  On an indecomposable Banach space, every idempotent has
finite-dimensional range or finite-dimensional kernel.  Thus a sum of two
idempotents is always a finite-rank perturbation of exactly one of
$0$, $\I$, or $2\I$.

Choose an indecomposable separable space without the compact approximation
property.  If sums of two idempotents converged strongly to
$\tfrac12\I$, one of the three scalar types would occur along an infinite
subsequence.  Removing that fixed scalar gives finite-rank operators converging
strongly to a nonzero scalar multiple of the identity.  Uniform boundedness then
turns them into bounded finite-rank approximants to $\I$, which converge uniformly
on compact sets.  This is the bounded approximation property and contradicts the
choice of the space.

\section*{Counterexample}

\begin{lemma}\label{lem:idempotent}
Let $X$ be an indecomposable infinite-dimensional Banach space.  For every
idempotent $P\in\Lop(X)$ there are $\varepsilon\in\{0,1\}$ and a finite-rank
operator $F$ such that
\[
P=\varepsilon\I+F.
\]
\end{lemma}

\begin{proof}
Since $P^2=P$, both $P[X]$ and $\ker P$ are closed and
\[
X=P[X]\oplus\ker P.
\]
Indecomposability forces one of these two summands to be finite-dimensional.
If $P[X]$ is finite-dimensional, take $\varepsilon=0$ and $F=P$.
If $\ker P$ is finite-dimensional, then $\I-P$ is finite-rank, so take
$\varepsilon=1$ and $F=P-\I$.
\end{proof}

\begin{lemma}\label{lem:bap}
Suppose finite-rank operators $R_j\in\Lop(X)$ satisfy
$R_jx\to x$ for every $x\in X$ and $\sup_j\|R_j\|<\infty$.  Then $X$ has
the bounded approximation property, and therefore the compact approximation
property.
\end{lemma}

\begin{proof}
Write $M=\sup_j\|R_j\|$.  Let $K\subset X$ be compact and $\eta>0$.
Choose a finite $\eta/(3(M+1))$-net $\{x_1,\ldots,x_m\}$ in $K$.
For all sufficiently large $j$,
$\max_i\|R_jx_i-x_i\|<\eta/3$.  If $x\in K$ and
$\|x-x_i\|<\eta/(3(M+1))$, then
\[
\|R_jx-x\|
 \le M\|x-x_i\|+\|R_jx_i-x_i\|+\|x_i-x\|<\eta.
\]
Hence $R_j\to\I$ uniformly on compact subsets, with uniformly bounded norms.
This is the bounded approximation property.  Since finite-rank operators are
compact, it implies the compact approximation property.
\end{proof}

\begin{theorem}\label{thm:main}
There exists a separable infinite-dimensional Banach space $X$ for which
$\tfrac12\I_X$ is not a strong-operator limit of sums of two idempotents.
Consequently Question 5.2.9 of arXiv:2307.03064 has a negative answer.

More generally, on the chosen space, $\lambda\I_X$ is not such a limit for
any scalar $\lambda\notin\{0,1,2\}$.
\end{theorem}

\begin{proof}
Allexandrov, Kutzarova, and Plichko prove that there is a closed subspace $X$
of the Gowers--Maurey hereditarily indecomposable space which fails the compact
approximation property~\cite{AKP}.  In particular, $X$ is separable,
infinite-dimensional, and indecomposable.

Fix $\lambda\notin\{0,1,2\}$ and suppose, toward a contradiction, that
idempotents $P_n,Q_n\in\Lop(X)$ satisfy
\[
P_n+Q_n\xrightarrow{\mathrm{sot}}\lambda\I.
\]
By Lemma~\ref{lem:idempotent}, write
\[
P_n=\varepsilon_n\I+F_n,\qquad
Q_n=\delta_n\I+G_n,
\]
where $\varepsilon_n,\delta_n\in\{0,1\}$ and $F_n,G_n$ are finite-rank.
Put $k_n=\varepsilon_n+\delta_n\in\{0,1,2\}$.  One value
$k\in\{0,1,2\}$ occurs along an infinite subsequence $(n_j)$.  Along this
subsequence,
\[
F_{n_j}+G_{n_j}\xrightarrow{\mathrm{sot}}(\lambda-k)\I.
\]
The uniform boundedness principle gives
$\sup_j\|F_{n_j}+G_{n_j}\|<\infty$.  Because $\lambda-k\ne0$, the
finite-rank operators
\[
R_j=(\lambda-k)^{-1}(F_{n_j}+G_{n_j})
\]
are uniformly bounded and converge strongly to $\I$.  Lemma~\ref{lem:bap}
then gives the compact approximation property for $X$, contradicting its
construction.  Taking $\lambda=\tfrac12$ gives the announced counterexample.
\end{proof}

\section*{Verification Notes}

No bound on the individual idempotents is assumed or needed.  Strong
convergence bounds the sums $P_n+Q_n$; after the scalar type is fixed on a
subsequence, it therefore bounds the finite-rank remainders used in the proof.
The changing scalar types cause no difficulty because there are only three of
them.  The distinction between pointwise convergence and uniform convergence
on compact sets is handled explicitly in Lemma~\ref{lem:bap}.

\section*{Novelty Check and Limitations}

The local run indexes were searched for arXiv:2307.03064, the source title,
the exact question wording, strong limits and sums of two idempotents, and the
bounded/compact approximation properties.  No matching packet or attempt was
found.  The surrounding source section was checked and does not answer
Question 5.2.9.

A bounded web search on August 26, 2026 using the exact question and close
strong-limit phrases found no later answer or this approximation-property
obstruction.  This was not a comprehensive MathSciNet or zbMATH search, so
novelty confidence is moderate.

The argument disproves the universal Banach-space statement.  It does not
classify the separable Banach spaces on which every operator has such an
approximation, nor does it weaken the HI/CAP construction to a familiar
classical space.

\section*{Human Review Recommendation}

Send to a Banach-space expert.  The proof itself is elementary once the
supporting existence proposition is accepted; the main review item is the
scalar-field scope of that standard HI construction.

\begin{thebibliography}{9}
\bibitem{Bourin}
J.-C. Bourin,
\emph{A Journey into Matrix Analysis},
arXiv:2307.03064, 2023.

\bibitem{AKP}
G. Allexandrov, D. Kutzarova, and A. Plichko,
\emph{A Separable Space with No Schauder Decomposition},
Proceedings of the American Mathematical Society 127 (1999), no.~9,
2805--2806,
\href{https://doi.org/10.1090/S0002-9939-99-05370-8}{doi:10.1090/S0002-9939-99-05370-8}.
\end{thebibliography}

\end{document}
